\documentclass[11pt]{letter}

\usepackage[margin=0.75in]{geometry}
\usepackage{setspace}
\usepackage{hyperref}
\usepackage{xcolor}
\usepackage{enumitem}
\usepackage{parskip}

\hypersetup{
    colorlinks=true,
    urlcolor=black,
    linkcolor=black
}

\setlength{\parindent}{0pt}
\setstretch{1.05}
\date{}
\begin{document}

\begin{letter}{}

\opening{Dear Hiring Committee,}

I saw the job post on \textbf{Panasonic Avionics careers site} and thought I'd be a good fit for it.

I am a Master's student in Computer Science at the University of Cincinnati, graduating December 2026, looking for a role as an \textbf{MTS II, Software} at \textbf{Panasonic Avionics Corporation}. With hands-on experience building backend automation, distributed infrastructure, and cloud platforms across Go, Python, and AWS, I am excited to contribute to the team building the real-time event logging and streaming systems behind Panasonic's inflight entertainment and connectivity products.

\textbf{Here are reasons that make me a great fit for the role:}

\begin{itemize}[leftmargin=*, itemsep=0.45em]
    \item \textbf{Backend systems in Go and distributed infrastructure:} At Siemens Industry Software, I developed a Go and Terratest framework catching 95\% of infrastructure misconfigurations pre-production, and automated Teamcenter X infrastructure workflows with Python and REST APIs, cutting deployment time by 40\% across AWS environments.

    \item \textbf{Real-time production systems and observability:} I executed 44 production environment upgrades in 4 months, diagnosing live infrastructure issues across Linux, EC2, and distributed databases using CloudWatch log correlation, and designed an Ansible-based Blue-Green deployment framework that reduced production upgrade downtime by 50\% while improving fault tolerance.

    \item \textbf{Cloud-native and containerized systems on AWS:} I migrated 40+ production environments to Kubernetes using ArgoCD, Helm, and GitLab CI/CD, and have hands-on experience with Docker, Terraform, and AWS services (EC2, S3, SQS, CloudWatch) for building scalable, resilient distributed systems.

    \item \textbf{AI-assisted monitoring and event analysis:} My current graduate research builds a real-time telemetry-analysis pipeline for automated event reproduction, achieving a 25.7x latency reduction over a published baseline, work that directly strengthens my ability to apply AI/ML to intelligent monitoring and anomaly detection.

    \item \textbf{Engineering quality and collaboration:} As Quality Coach and SonarQube Specialist at Siemens, I increased automated code coverage from 35\% to 78\% and mentored 5 software engineers through code review and structured onboarding.
\end{itemize}

I am available for full-time work after my graduation in December 2026. As a Master's student at the University of Cincinnati and someone who values professionalism and growth, I am enthusiastic about contributing to Panasonic Avionics and helping engineer the real-time systems that bring entertainment and connectivity to passengers worldwide.

Please find my CV attached below. I look forward to hearing from you. I would be excited to help build the technology that keeps millions of travelers connected and entertained at altitude.

{\setlength{\parskip}{0pt}
Regards,\\
Tanmay Bhuskute,\\
tanmay.v.bhuskute@gmail.com $\mid$ +1 (732) 858-3027,\\
Cincinnati, OH.
}

\end{letter}

\end{document}
